\documentclass[11pt, letterpaper]{article}

\usepackage[
    ignoreheadfoot,
    top=2 cm,
    bottom=2 cm,
    left=2 cm,
    right=2 cm
]{geometry}
\usepackage[dvipsnames]{xcolor}
\definecolor{primaryColor}{RGB}{0, 0, 0}
\usepackage{parskip}
\usepackage{charter}
\usepackage[
    pdftitle={Sreeniketh Aathreya Cover Letter},
    pdfauthor={Sreeniketh Aathreya},
    colorlinks=true,
    urlcolor=primaryColor
]{hyperref}

\pagestyle{empty}
\setlength{\parindent}{0pt}
\raggedright

\begin{document}

{\fontsize{18 pt}{20 pt}\selectfont\bfseries Sreeniketh Aathreya}\\[4pt]
Software Engineer / Solutions Architect\\[6pt]
\href{https://sreenikethaathreya.com}{sreenikethaathreya.com}
~\textbullet~
\href{mailto:sreenikethaathreya@gmail.com}{sreenikethaathreya@gmail.com}
~\textbullet~
\href{tel:+1-571-279-7924}{+1-571-279-7924}
~\textbullet~
\href{https://www.linkedin.com/in/sreeniketh-aathreya-16800b122/}{LinkedIn}

\vspace{12pt}
\hrule
\vspace{12pt}

September 28, 2026

\vspace{8pt}

Hiring Team\\
Cvent

\vspace{12pt}

Dear Hiring Team,

Cvent's Technical Architect - AI role is the work I already do: folding large language models, agents, and retrieval-augmented generation into enterprise SaaS so products stay scalable, reliable, and usable. What draws me to Cvent is the platform itself---software that brings millions of people together at events---and a culture that treats AI as part of how the company builds, not a side experiment. I want to help productize those capabilities across your event marketing and management stack with the architectural quality and paved-road discipline this role describes.

As Solutions Architect and AI Labs Lead at HatchWorks AI, I partner with the CTO on technical direction for 20+ client engagements, translating business goals into AI-native architectures, staffing estimates, and delivery-ready proposals. I was a major contributor to GenDD, an organization-wide AI delivery framework covering 150+ software-development tasks; adoption reclaimed 347 AI hours per month and produced a 75\% productivity gain with under a 3\% defect rate---the reuse, training, and developer-productivity outcome this posting calls for. I design agentic workflows with LangChain, embeddings, and foundation models, and I run build-versus-buy evaluations across AWS, GCP, and Azure so teams inherit a coherent cloud strategy instead of a pile of tools.

The same pattern shows up in how I ship and how I communicate. At i-Link I helped architect an event-driven cloud-native platform that processes millions of API calls daily at 99.9\% uptime, and I mentored junior engineers on microservices and Agile. At R3 I closed a \$4 million enterprise deal as a Pre-Sales Engineer; at Akamai I influenced \$20M+ in platform deals through POCs and sessions with C-level and technical audiences. I implement RAG and access control on Amazon Bedrock, write about permission-aware GenAI and RAGOps on Medium and LinkedIn Pulse, and demo agent systems on YouTube so practices can be reused. I build in TypeScript, Python, and Java, and I ship AI features on AWS (Lambda, CloudFormation, S3, CloudWatch, Bedrock) with CI/CD automation and persistence layers such as PostgreSQL and DynamoDB.

I am based in Fairfax / Reston, Virginia and ready to talk through how Cvent can integrate LLMs, agents, and RAG into existing products without sacrificing performance or judgment. I would welcome a conversation with the hiring team on this Technical Architect - AI role.

Sincerely,\\[12pt]
Sreeniketh Aathreya

\end{document}
